% Generated draft: requires mathematical and source review.
\documentclass[11pt]{article}
\usepackage[margin=25mm]{geometry}
\usepackage{fontspec}
\setmainfont{Latin Modern Roman}
\usepackage{amsmath,amssymb,mathtools,unicode-math}
\setmathfont{Latin Modern Math}
\usepackage{enumitem,needspace,xurl,hyperref}
\setlength{\parindent}{0pt}
\setlength{\parskip}{5pt}
\setlength{\emergencystretch}{3em}
\newcommand{\N}{\mathbb{N}}
\newcommand{\Z}{\mathbb{Z}}
\newcommand{\Q}{\mathbb{Q}}
\newcommand{\R}{\mathbb{R}}
\newcommand{\C}{\mathbb{C}}
\newcommand{\F}{\mathbb{F}}
\newcommand{\A}{\mathbb{A}}
\newcommand{\PP}{\mathbb{P}}
\newcommand{\E}{\mathbb{E}}
\newcommand{\eps}{\varepsilon}
\newcommand{\dd}{\,\mathrm{d}}
\begin{document}
\section*{UnsolvedMath ID: 11000296}
\textbf{Source problem code:} catalogue entry
\subsection{Definitions and mathematical statement}
Question 5.4 asks whether the image of the second Morita class in the source's degree-\(8\) rational group for \(\mathrm{GL}(6,\Z)\) is nonzero. The Morita class, source group and inducing map are not defined in the extract, and the flattened notation \(H8\) does not preserve a subscript/superscript distinction.

A natural conventional interpretation uses the abelianization homomorphism
\[
 a:\operatorname{Out}(F_6)\longrightarrow\mathrm{GL}(6,\Z),
\]
induced from the action on \(F_6/[F_6,F_6]\cong\Z^6\). If the source's second Morita class is the corresponding \(\mu_2\in H_8(\operatorname{Out}(F_6);\Q)\), the question is
\[
 a_*(\mu_2)\ne0\quad\text{in }H_8(\mathrm{GL}(6,\Z);\Q)\ ?
\]
This homological reading is explicitly conditional, not a recovered construction of \(\mu_2\). Reading the target as cohomology \(H^8\) would require a different specified map, because a group homomorphism induces cohomology pullback in the opposite direction. The exact class and variance must therefore be checked against the source. Rank \(6\), degree \(8\) and rational coefficients are fixed; no stable-rank or integral substitute is made. The source points to [81] for further discussion of \(\operatorname{Aut}(F_n)\) and \(\operatorname{Out}(F_n)\) (co)homology but supplies no definition or computation here.
\subsection{Short English statement}
Is the specified second Morita class nonzero after passage to \(\mathrm{GL}(6,\Z)\) in rational degree eight? The natural homological image is displayed conditionally, since the class definition, map and lost homology/cohomology grading need source confirmation.
\subsection{Sources}
\begin{itemize}
\item Local UnsolvedMath registry, permanent ID \(11000296\), exact source code \texttt{catalogue entry}, Question 5.4, PDF page 337. No individual proposer or proposed year is supplied.
\item Farb, edited volume \emph{Problems on Mapping Class Groups and Related Topics}: \url{https://www.math.uchicago.edu/~farb/papers/mcgbook.pdf}. The volume has 329 explicitly typeset numbered items across chapters, repeated local numbering distinguished by PDF page.
\item Internal reference [81] for further free-group automorphism/outer-automorphism cohomology discussion; full entry absent. The second Morita class has no separate construction reference in the record.
\item Extraction pdf-font-structure, accessed 2026-07-29. Source evidence NEEDS\_REVIEW presents an open problem and calls for a current primary-literature audit. Supplied OPEN-TRIAGE summary mentions Morita classes and stable/unstable linear-group cohomology but reports no verified decisive computation.
\item Dataset attribution: Ulam AI / UnsolvedMath Contributors, CC BY 4.0. Source rights NEEDS\_REVIEW: publicly accessible author/publisher copy, redistribution terms requiring review.
\end{itemize}
\end{document}
